\scanpage{9}
% archive foliation: 8
pour $j$ \inserted{$\Gamma$-}gagnantes, \deleted{pour tt} $\Sigma$ est la plus grande, i.e. $\Sigma\supset\Sigma'$, i.e. $\forall z\in C(j)$, on a $\Sigma(z)\supset\Sigma'(z)$. Bien évident (puisque $\Sigma'(z)\subset R(z)$) \deleted{\unsure{si $\Sigma$}} dans le cas où $\Sigma(z)=R(z)$, donc quand \deleted{\unsure{on n'a pas simultanément}}\inserted{$z\notin\Gamma$} \deleted{$\exists x\in\Gamma$, avec $z\unsure{\succ|\geq}x$, et $z\in G(j)$.} On peut supposer \deleted{\unsure{ces conditions satisfaites}}\inserted{$z\in\Gamma$}, donc $\Sigma(z)=R(z)\cap G(j)$. Il suffit donc de prouver que $\Sigma'(z)\subset G(j)$, \deleted{\unsure{ce}} ce qui résulte du \uuline{lemme} après la prop.

\deleted{Stratégie}

$4^\circ)$ \uline{Stratégie acceptable}, \uline{positions acceptables}.

$J'$-Stratégie (si $J'\subset J$) \deleted{\unsure{Se}} \unsure{Posant|Posons}
\[
C(J')=\bigcup_{j\in J'}C(j),
\]
se définit comme une $j$-stratégie, en y remplaçant $C(j)$ par $C(J')$.
Donc les $j$-stratégies deviennent les $\{j\}$-stratégies.
\margin{\unsure{pour $x\in C$|pour celle}}

Une \deleted{\unsure{telle}} $J'$-stratégie \inserted{$\Sigma'$} est dite \uline{acceptable} (\deleted{pour}\inserted{\uline{contre}} un $j_0\in J$ \inserted{et $x\in C$}) si pour $\forall$ partie \inserted{$(x_0,\ldots,x_N)$} comp. avec $\Sigma$, \deleted{\unsure{avec}}\inserted{partant de $x$}\deleted{\unsure{terminant par}}, et terminée, on a $x_N\notin G(j_0)$
\inserted{i.e. $x_N\notin G_0(j_0)$.}
